\begin{example}[Compatible pair but several points]\label{ex:vi23-a01}
For a field $k$ and finite extensions $K,L$, points over the unique pair in $\Spec K\times_k\Spec L$ are the primes of $K\otimes_kL$, which can be more than one.
\end{example}

\begin{example}[Separable splitting phenomenon]\label{ex:vi23-a02}
If $L/k$ is finite separable and $\Omega$ contains all $k$-embeddings of $L$, then $L\otimes_k\Omega$ can split as a product of copies of $\Omega$, yielding several points above one set-theoretic pair.
\end{example}

\begin{example}[Inseparable tensor product]\label{ex:vi23-a03}
In characteristic $p$, tensor products of purely inseparable extensions may be nonreduced, showing that a fiber product can acquire nilpotents even when both factors are fields.
\end{example}

\begin{example}[Equalizers from diagonals]\label{ex:vi23-a04}
For $f,g:T\rightrightarrows X$ over $S$, the equalizer condition $f=g$ is imposed by pulling the map $(f,g):T\to X\times_SX$ back along $\Delta_{X/S}$.
\end{example}

\begin{example}[Intersections are scheme-theoretic]\label{ex:vi23-a05}
When two closed subschemes of $X$ meet, their fiber product remembers the sum of defining ideals, including nilpotent structure invisible to the set-theoretic intersection.
\end{example}

\begin{example}[Localization commutes with tensor product]\label{ex:vi23-a06}
For $f\in A$, $(B\otimes_AC)\otimes_AA_f\cong B_f\otimes_{A_f}C_f$, explaining compatibility of affine pullbacks with restriction to principal opens.
\end{example}

\begin{example}[Descent of the affine construction to overlaps]\label{ex:vi23-a07}
On overlapping affine charts, the two tensor-product descriptions localize to canonically isomorphic rings. Uniqueness supplies the cocycle needed for gluing.
\end{example}

\begin{example}[Graph and closedness preview]\label{ex:vi23-a08}
The graph is always a pullback of the diagonal. Therefore any future property of the diagonal stable under pullback transfers to graphs; this observation underlies separatedness theory.
\end{example}

\begin{example}[Products and final objects]\label{ex:vi23-a09}
The existence of the final object $\Spec\mathbb Z$ converts binary products in schemes into fiber products over an absolute base.
\end{example}

\begin{example}[Iterated compatibility]\label{ex:vi23-a10}
An iterated fiber product represents triples of maps with one common image in the base, explaining canonical associativity without choosing coordinates.
\end{example}

\begin{example}[Functor-of-points viewpoint]\label{ex:vi23-a11}
For every $T$, $(X\times_SY)(T)$ is the set-theoretic fiber product $X(T)\times_{S(T)}Y(T)$. This does not mean underlying topological points form a naive fiber product.
\end{example}

\begin{example}[Tensor product zero]\label{ex:vi23-a12}
It is possible for a tensor product over a base ring to be the zero ring; then the affine fiber product is empty. Compatibility over a base can therefore eliminate all geometric points.
\end{example}

\begin{example}[Residue fields at a product point]\label{ex:vi23-a13}
A product point determines maps from $\kappa(x)$ and $\kappa(y)$ into its residue field compatible over the base. This is the field-theoretic refinement of the pair $(x,y)$.
\end{example}

\begin{example}[Locality of existence]\label{ex:vi23-a14}
The global existence theorem is a model local-to-global argument: solve the universal problem on affine charts, verify compatibility on overlaps, then glue the representing objects.
\end{example}
